% HA-Expert: A Domain-Adapted Language Model and Evidence-Checked
% Benchmark for Hyaluronan Research
% arXiv preprint manuscript (cs.CL primary; cross-list cs.LG, q-bio.BM).
% Sources: notes/arxiv-draft-v1.md; eval/results/s6-full-report.md;
% eval/verify_summary_v3.md; notes/PROJECT_PLAN.md; model card
% (hf-staging/model-card-sft.md). All reported numbers are taken from
% those records.
% Pre-submission notes: (1) an arXiv endorsement code is still required
% before submission; (2) the reference list is the draft's selected list
% and may be completed at submission time.
\documentclass[11pt]{article}

\usepackage[letterpaper,margin=1in]{geometry}
\usepackage{booktabs}
\usepackage{array}
\usepackage{amsmath}
\usepackage{url}
\usepackage{hyperref}
\hypersetup{colorlinks=true, linkcolor=blue, citecolor=blue, urlcolor=blue}

\title{HA-Expert: A Domain-Adapted Language Model and Evidence-Checked
Benchmark for Hyaluronan Research}

\author{Kunli Liu\\
Independent Researcher, San Francisco Bay Area, USA\\
\href{mailto:liukunli001@gmail.com}{liukunli001@gmail.com}}

\date{}

\begin{document}
\maketitle

\begin{abstract}
Specialized scientific fields accumulate knowledge in large, dispersed
literatures that general-purpose language models handle unreliably: answers
are fluent but often wrong on quantitative details, and confident even when
the premise of a question is false. We study one such field --- hyaluronan
(hyaluronic acid, HA), a glycosaminoglycan whose research spans
glycochemistry, biomaterials, drug delivery, and clinical medicine --- and
ask how far a single researcher with a modest compute budget can go in
building a \emph{trustworthy} domain model. We release (i) a cleaned corpus
of 12,306 full-text HA papers ($\sim$126M tokens); (ii) HA-Expert, a
Qwen3-8B-Base model adapted in two stages --- domain-adaptive continued
pretraining (DAPT) on 60M tokens of the corpus, followed by full-parameter
supervised fine-tuning (SFT) on 12,003 question--answer pairs generated by
the DAPT model itself, at zero API cost --- and optionally grounded at
inference time by retrieval over a 636,959-passage FAISS index of the same
corpus; and (iii) HA-Expert-Bench, 240 questions whose reference answers
were verified against retrieved source passages by evidence-restricted LLM
verifiers (212 of 240 keys confirmed or corrected through this process),
including 15 adversarial ``trap'' questions with false premises. In a
four-cell ablation, SFT doubles the base model (0.300 $\rightarrow$ 0.601);
retrieval is the largest single lever (base+RAG 0.710; traps
0.267 $\rightarrow$ 0.867); and, notably, SFT and RAG do \textbf{not}
compose additively (0.662): a model that never saw retrieved context in
training collapses into citation-marker loops on 5 of 240
retrieval-augmented responses, and inherits its teacher's repetition habit
through self-distillation. Total cloud compute for the full pipeline was
under USD 50. All code, data, weights, and the benchmark are publicly
released.
\end{abstract}

\section{Introduction}
\label{sec:intro}

Hyaluronan is a linear glycosaminoglycan of repeating
$\beta$-1,4-glucuronic-acid / $\beta$-1,3-$N$-acetylglucosamine
disaccharides, central to extracellular-matrix biology, joint lubrication,
wound healing, ophthalmic and dermal formulations, and an expanding
drug-delivery literature. The knowledge a practitioner needs ---
molecular-weight-dependent biology, receptor interactions (CD44, RHAMM,
LYVE-1), chemical modification routes, production methods, clinical
dosing --- is spread across thousands of papers in chemistry, materials,
and medical journals. No single textbook tracks it, and the literature
grows faster than any individual's reading.

General-purpose LLMs are an attractive interface to such a literature.
However, in our initial probes their failure modes in this domain were
systematic rather than occasional, and they were failures of three
distinct kinds. First, a failure of \emph{behavior}: the base model we
started from frequently continued the prompt with invented follow-up
questions instead of answering the question asked. Second, a failure of
\emph{epistemic stance}: the model assented to false premises --- for
example, that HA is a protein polymer, or that HA solutions are Newtonian
fluids --- rather than correcting them. Third, a failure of
\emph{verifiability}: quantitative claims were produced fluently, with no
checkable source a domain expert could audit. A system that fails in these
three ways is not merely inaccurate; it is untrustworthy in exactly the
situations --- unfamiliar quantitative detail, adversarially false
premises --- where a specialist tool is most needed.

This paper is a case study in fixing those failure modes with the standard
toolkit --- continued pretraining, instruction tuning, retrieval --- under
three constraints that we believe are typical for individual researchers
and small labs: (1) a single rented GPU and a total compute budget under
USD 50; (2) no paid teacher model --- all supervision is self-generated
and then filtered; and (3) evaluation that a domain expert can trust,
which for us meant reference answers that are \emph{checked against source
text}, not merely model-written. Each constraint shapes the design.
Constraint (1) fixes the model scale at 8B parameters and the training
schedule at hours, not weeks. Constraint (2) turns the pipeline into a
self-distillation exercise, whose characteristic risks --- the teacher's
errors and habits propagating into the student --- we measure rather than
assume away. Constraint (3) makes benchmark construction a research
artifact in its own right: every answer key in HA-Expert-Bench carries a
verification status obtained by checking it against retrieved source
passages, and we report the full accounting, including the keys that
verification \emph{changed} and the keys that remain uncertain.

The resulting system, HA-Expert, couples a two-stage adapted model
(Section~\ref{sec:methods}) with an optional retrieval mode over an index
of the same corpus it was trained on. On HA-Expert-Bench
(Section~\ref{sec:bench}), fine-tuning doubles the base model's score,
retrieval applied to the \emph{untrained} base model is worth still more,
and the combination of the two --- the configuration a practitioner would
most naturally deploy --- underperforms retrieval alone, for a reason that
output inspection makes precise (Section~\ref{sec:analysis}). We regard
that negative interaction, and the evidence-checked evaluation protocol
that makes it measurable, as the paper's most transferable findings.

Contributions:
\begin{itemize}
\item \textbf{Corpus.} A documented, decontaminated corpus of 12,306
open-access HA papers ($\sim$126M tokens), with 30 review papers held out
from every stage --- training, SFT data generation, and the retrieval
index (Section~\ref{sec:corpus}).
\item \textbf{Model.} HA-Expert (Qwen3-8B-Base + DAPT + SFT), adapted
entirely with self-generated supervision, plus a retrieval mode with
inline citations; weights, a DAPT checkpoint, and a public demo are
released (Section~\ref{sec:methods}).
\item \textbf{Benchmark.} HA-Expert-Bench: 240 questions with an
evidence-verification pipeline for the answer keys (two rounds:
keyword-evidence, then vector-evidence re-verification), and an explicit
accounting of how many keys are confirmed, corrected, or still uncertain
(Section~\ref{sec:bench}).
\item \textbf{Ablations and failure analysis.} A four-cell study
(base / +SFT / +RAG / +both) with per-type scores, a trap-by-trap table,
and a documented interference effect between SFT and RAG with a concrete
remediation path (Sections~\ref{sec:experiments}
and~\ref{sec:analysis}).
\end{itemize}

\section{Related Work}
\label{sec:related}

Our approach deliberately assembles standard components rather than new
algorithms; the interest is in how they interact under the constraints of
Section~\ref{sec:intro}, and in the evaluation protocol that makes those
interactions measurable.

\textbf{Domain-adaptive pretraining and fine-tuning.} Continued
pretraining of a general base model on domain text, followed by
full-parameter supervised fine-tuning, is the established route to domain
language models. We instantiate it with Qwen3-8B-Base as the base model
and the corresponding Qwen3-Embedding-0.6B as the retrieval
embedder~\cite{qwen3}, using the Transformers and TRL
libraries~\cite{transformers,trl}. Two details of the instantiation are
worth noting because they recur in the analysis: SFT used NEFTune noisy
embeddings~\cite{neftune}, and the final phase of DAPT mixed in
general-domain replay text from FineWeb-Edu~\cite{fineweb} against
catastrophic forgetting.

\textbf{Self-generated supervision.} Where a frontier teacher is
unaffordable, a model can generate its own instruction data; the known
hazard is that the generator's errors and stylistic habits are distilled
into the student unchecked. Our mitigation is a grounding filter ---
candidate answers are retained only if they are attested by the corpus
passage they were generated from (Section~\ref{sec:sft-data}) --- and our
evaluation is designed, in part, to detect the failure modes that survive
the filter. The teacher was served locally with vLLM~\cite{vllm}.

\textbf{Retrieval grounding.} Dense retrieval over the domain corpus,
here FAISS inner-product search~\cite{faiss}, is the standard remedy for
parametric-knowledge failures: the model is shown the relevant passages
and asked to cite them. Our results complicate the standard story in one
specific way --- a model fine-tuned without ever seeing retrieved context
can be made \emph{worse} by adding it --- which we document with counts,
per-question verdicts, and verbatim case studies
(Section~\ref{sec:analysis}).

\textbf{Benchmark keys and LLM judges.} Domain benchmarks are commonly
shipped with model-written answer keys that are never re-checked, and
with LLM judges whose verdicts inherit whatever errors the keys contain.
Our contribution relative to that practice is procedural: an
evidence-restricted verification pipeline for the keys themselves, in
which verifiers may only confirm, correct, or abstain
(``uncertain'') on the basis of retrieved source passages, with
per-question verification status released alongside the scores
(Section~\ref{sec:bench}). Verification changed 22 of our 240 keys; we
treat that figure as an estimate of the silent error rate that unchecked
keys would otherwise inject into every reported score.

\section{Corpus}
\label{sec:corpus}

\textbf{Harvest.} We harvested records from PubMed Central's open-access
subset using hyaluronan-related queries, obtaining 16,642 records, of
which 12,945 had retrievable full text in JATS XML.

\textbf{Cleaning.} Full texts were converted from JATS XML to plain text.
Cleaning removed navigation and boilerplate material; reference lists were
retained as text, since citation context is part of the literature's
content. Exact- and near-duplicate documents were dropped, and the
remainder was language- and length-filtered. The final corpus contains
\textbf{12,306 papers, $\approx$126M tokens}. The corpus and its manifests
are released (GitHub Release v0.1-data).

\textbf{Holdout.} Thirty review papers were set aside \emph{before} any
processing and are excluded from every downstream stage: the DAPT sets,
the passages used to generate SFT data, and the retrieval index.
Benchmark questions derived from these papers therefore measure
generalization to literature the system has never trained on or retrieved
from, rather than recall.

\textbf{DAPT sets.} The continued-pretraining data comes in two parts.
The \emph{core set} contains 4,952 papers (60.0M tokens). The
\emph{anneal set} contains 498 papers (5.7M tokens) of high-quality
reviews and methods papers, placed at the end of training and mixed with
12\% general-domain replay text (FineWeb-Edu) against catastrophic
forgetting: the intent is a short final phase on the corpus's most
reliable text, with replay to keep the model's general abilities from
drifting as the domain phase ends.

\section{Methods}
\label{sec:methods}

The pipeline has three stages --- domain-adaptive pretraining, supervised
fine-tuning on self-generated data, and retrieval-augmented inference ---
followed by evaluation on HA-Expert-Bench (Section~\ref{sec:bench}). We
specify each stage at the level needed to replicate it. All training ran
on a single rented GPU under constraint (1) of Section~\ref{sec:intro}.

\subsection{Domain-adaptive pretraining (DAPT)}
\label{sec:dapt}

Starting from Qwen3-8B-Base~\cite{qwen3}, we ran full-parameter continued
pretraining on the DAPT sets of Section~\ref{sec:corpus}. Documents were
chunked at document boundaries to a sequence length of 4,096 tokens
(legacy block packing) and processed with SDPA attention. The effective
batch was 8 sequences $\times$ gradient accumulation 16, i.e.\ 128
sequences or 524,288 tokens per step. The learning rate was
$1\times10^{-5}$ with cosine decay. Training ran for 212 steps
($\approx$111M tokens processed in total) in 5\,h 34\,min on a single
RTX PRO 6000 (96\,GB), reaching a final training loss of $\approx$1.55.

A development check after this stage set the agenda for the rest of the
pipeline. On a 63-question development split, the DAPT model tied the base
model on the auto-graded items (5/6), even as it repaired knowledge traps
the base model missed. In other words, DAPT changed what the model
\emph{knows} without changing how it \emph{behaves}: left to itself, it
still completed prompts rather than answering them. That division of labor
--- DAPT for knowledge, a second stage for behavior --- motivated the SFT
stage, and the full-bench results of Section~\ref{sec:experiments} bear
it out.

\subsection{Self-generated SFT data}
\label{sec:sft-data}

Paying for a frontier teacher was out of budget, so the DAPT model served
as its own teacher. The generation protocol was:
\begin{enumerate}
\item Serve the DAPT model locally with vLLM~\cite{vllm} (zero API cost).
\item Sample passages from the corpus, excluding the 30 holdout papers of
Section~\ref{sec:corpus}.
\item Generate a question--answer pair conditioned on each passage, so
that every candidate answer is anchored to a specific source text.
\item Retain a pair only if the answer is \emph{grounded} in its source
passage, operationalized as a content-word overlap of $\geq 0.55$ between
the answer and the passage; deduplicate the survivors.
\end{enumerate}
Two generation rounds produced 5,403 and 6,600 retained pairs,
respectively, for \textbf{12,003 examples} in total. The grounding filter
is a check of topical attestation, not of truth: it guarantees that the
answer's content words appear in a real passage, not that the answer
states the passage's content correctly. We return to this distinction in
Section~\ref{sec:limitations}; it is one reason the benchmark's keys, by
contrast, were verified by a separate evidence-restricted process
(Section~\ref{sec:bench}).

\subsection{Supervised fine-tuning}
\label{sec:sft-train}

SFT was full-parameter, in prompt/completion format with loss computed on
the completion only, at learning rate $1\times10^{-5}$ for 2 epochs
(752 steps), with NEFTune noise~\cite{neftune} 5.0 and an 8-bit Adam
optimizer. Training finished at train loss 1.214 in $\approx$19 minutes. The
qualitative change was immediate and matches the quantitative one: the
base model's signature pathology --- inventing follow-up questions instead
of answering --- disappeared from the fine-tuned model's outputs.

\subsection{Retrieval-augmented inference (RAG)}
\label{sec:rag}

For the retrieval mode, all 12,306 papers were chunked by merging
paragraphs to $\approx$900 characters, with a 150-character overlap on
long paragraphs, and embedded with Qwen3-Embedding-0.6B
(1,024-dimensional), giving a FAISS IndexFlatIP~\cite{faiss} index of
\textbf{636,959 passages}. Holdout papers are absent from the index, as
they are from every training stage. At inference, the top-4 passages for
a question are prepended to it with citation markers [1]--[4], and the
model is instructed to cite them in its answer.

One decoding detail matters for interpreting the results. The public demo
decodes with a repetition penalty of 1.05, which largely suppresses the
repetition failure mode analyzed in Section~\ref{sec:analysis}. The
benchmark evaluation of Section~\ref{sec:experiments} deliberately uses
plain greedy decoding with no repetition penalty, so that decoding aids do
not mask the behavior being measured.

\section{HA-Expert-Bench}
\label{sec:bench}

\subsection{Composition}

HA-Expert-Bench contains 240 questions spanning seven content categories:
basics; history and industry; synthesis and chemistry; modification and
derivatives; molecular-weight biology and receptors; nano-delivery; and
experimental and clinical topics. Questions take four formats
(Table~\ref{tab:benchcomp}). \emph{Trap} questions embed a false premise
that a knowledgeable expert should reject rather than answer on its own
terms --- e.g., ``HA solutions are simple Newtonian fluids, so their
viscosity is a single constant, correct?'' The trap format exists because
premise rejection was one of the three failure modes identified in
Section~\ref{sec:intro}, and it is invisible to benchmarks that only ask
well-formed questions.

\begin{table}[t]
\centering
\begin{tabular}{lcll}
\toprule
Format & $n$ & Grading & Notes \\
\midrule
Open & 195 & LLM judge & free-form answers \\
Quantitative & 20 & automatic & numeric tolerance \\
Trap & 15 & LLM judge & false premise to reject \\
Multiple-choice & 10 & automatic & exact option letter \\
\midrule
Total & 240 & & \\
\bottomrule
\end{tabular}
\caption{HA-Expert-Bench composition. Judge-graded items use the
correct / partial / wrong rubric of Section~\ref{sec:grading}, scored
1.0 / 0.5 / 0.0 against the evidence-checked reference key.}
\label{tab:benchcomp}
\end{table}

\subsection{Evidence-checked answer keys}
\label{sec:keys}

Draft keys were not trusted as written. Each key passed through an
evidence-restricted verification protocol in which LLM verifiers received
an evidence pack of source passages and could take exactly three actions:
\emph{confirm} the key, \emph{correct} it while quoting the decisive
passage, or mark it \emph{uncertain}. Verifiers were not permitted to
assert facts from outside the pack.

\textbf{Round one (keyword evidence).} Evidence packs were assembled by
tracing each question to its source paper (38 questions) or by keyword
retrieval over the corpus (202 questions). Outcome: 118 keys verified,
19 corrected, 103 marked uncertain.

\textbf{Round two (vector evidence).} The 103 uncertain keys were
re-verified against evidence packs of the top-6 passages retrieved from
the vector index of Section~\ref{sec:rag}, under the same three-action
protocol (six parallel verifier batches; 103 of 103 judged). Outcome:
72 verified, 3 corrected, 28 still uncertain.

\textbf{Released key set (v3).} Across both rounds, 190 keys are verified
and 22 corrected, so \textbf{212 of 240 keys (88.3\%) are
evidence-settled}; the remaining 28 (11.7\%) are released flagged as
uncertain drafts. Only three keys changed in round two (B22 and F07,
quantitative; G14, trap). Because scores had already been computed
against the v2 keys, we re-scored these three items explicitly under the
v3 keys: the auto-grades for B22 and F07 are identical in all four cells,
and G14's verdicts are unchanged on re-judging. All scores reported in
Section~\ref{sec:experiments} therefore stand under the released v3 keys.
Every per-question verification status is published with the benchmark.

\subsection{Grading protocol}
\label{sec:grading}

Multiple-choice items are graded by exact option letter and quantitative
items by numeric tolerance; both are fully automatic. Open and trap items
are graded by an LLM judge shown the question, the evidence-checked
reference answer, and the model's response, and required to return a
verdict of correct, partial, or wrong, scored 1.0, 0.5, and 0.0
respectively. The judge was the SFT model itself, served locally --- a
deliberate \$0 choice whose bias implications we treat as a limitation
(Section~\ref{sec:limitations}). Two features of the protocol mitigate,
without eliminating, that bias: the judge is anchored to a fixed verified
reference key rather than expressing a free pairwise preference, and the
30 automatic items provide a judge-free control
(Section~\ref{sec:experiments}). Judge outputs that failed to parse were
re-judged in a retry pass, which moved overall scores by at most 0.004;
the residual ungraded items are accounted for in
Section~\ref{sec:experiments}. All responses were generated with greedy
decoding (temperature 0) and a maximum of 512 tokens.

\section{Experiments and Results}
\label{sec:experiments}

\subsection{Protocol}

We evaluate four cells under a single protocol: the base model
(Qwen3-8B-Base); the adapted model (DAPT+SFT); and each of the two with
retrieval (base+RAG, DAPT+SFT+RAG). Responses for all 240 questions were
generated with greedy decoding (temperature 0, maximum 512 tokens) on a
single rented RTX PRO 6000 (96\,GB), and graded as specified in
Section~\ref{sec:grading}. Scores are means over graded items. A small
number of items remained ungraded because the judge's output was
unparseable even after the retry pass (judge JSON failures): 13 for the
base model, 23 for DAPT+SFT, 2 for base+RAG, and 21 for DAPT+SFT+RAG,
leaving $n = 227$, $217$, $238$, and $219$ graded responses per cell,
respectively. Following the judge's known limits, we treat cross-cell gaps
below $\approx$0.02 as noise throughout.

\begin{table}[t]
\centering
\begin{tabular}{lccccc}
\toprule
Cell & Overall & MCQ (10) & Quant.\ (20) & Open (195) & Trap (15) \\
\midrule
Qwen3-8B-Base        & 0.300 & 0.800 & 0.250 & 0.282 & 0.267 \\
+ DAPT \& SFT        & 0.601 & \textbf{0.900} & 0.300 & 0.631 & 0.500 \\
+ RAG                & \textbf{0.710} & \textbf{0.900} & 0.450 & \textbf{0.719} & \textbf{0.867} \\
+ DAPT \& SFT \& RAG & 0.662 & 0.700 & \textbf{0.550} & 0.676 & 0.633 \\
\bottomrule
\end{tabular}
\caption{HA-Expert-Bench results (240 questions; means over graded items,
$n$ per cell as in Section~\ref{sec:experiments}). Bold marks the best
score in each column. Three facts organize the table: fine-tuning doubles
the base model; retrieval on the \emph{base} model is worth more than the
full training pipeline without it; and the combination underperforms
retrieval alone on open and trap questions while posting the best
quantitative score. On the 30 auto-graded items alone (no judge involved):
0.433 / 0.500 / 0.600 / 0.600 in the same cell order.}
\label{tab:main}
\end{table}

\subsection{Fine-tuning doubles the base model}

DAPT+SFT scores 0.601 overall against the base model's 0.300. The gain is
largest exactly where the base model's deficit was behavioral: open
questions rise from 0.282 to 0.631, because the fine-tuned model answers
the question asked instead of continuing the prompt with invented
follow-ups. This confirms, at full-bench scale, the division of labor
observed in the development check of Section~\ref{sec:dapt}: DAPT
supplies domain knowledge, and SFT converts a text completer into a
system that responds to questions. Trap accuracy also improves
(0.267 $\rightarrow$ 0.500) but remains far from reliable --- a
parametric model, however well adapted, still assents to half of the
false premises it is offered.

\subsection{Retrieval is the largest single lever}

Adding retrieval to the \emph{base} model --- no training at all ---
scores 0.710, beating the full training pipeline without retrieval
(0.601). The effect is dramatic exactly where parametric knowledge is
least reliable: trap accuracy rises from 0.267 to 0.867. Both
long-standing trap failures of the parametric models (G01, ``HA is a
protein polymer''; G13, ``HA solutions are Newtonian fluids'') are
corrected by both retrieval cells. The pattern has a simple reading: a
false premise is hard to reject from memory, because rejection requires
confidence in a fact the model is unsure of; it is much easier to reject
when the contradicting passage is in the prompt. Retrieval converts an
epistemic problem into a reading problem, and reading is what base models
already do well.

\subsection{SFT and RAG interfere}

The combined cell, DAPT+SFT+RAG (0.662), trails base+RAG overall despite
dominating it without retrieval. The deficit is concentrated in the two
judge-graded formats: open (0.676 vs.\ 0.719) and trap (0.633
vs.\ 0.867). It is not uniform: on quantitative questions the combined
cell posts the best score of any cell (0.550, vs.\ 0.450 for base+RAG),
as if trained domain numeracy composes with retrieved figures when the
answer is a number to be extracted or computed, while the interference
appears where the model must produce grounded prose. A judge-free control
supports localizing the effect this way: on the 30 auto-graded items,
DAPT+SFT+RAG \emph{ties} base+RAG (0.600 each), so the headline gap lives
in the judged formats, not in the objective subset. Section~\ref{sec:analysis}
identifies the mechanism from the outputs themselves.

\section{Analysis}
\label{sec:analysis}

\subsection{Mechanism 1: a train/test format mismatch}

The SFT data of Section~\ref{sec:sft-data} consists exclusively of bare
question$\rightarrow$answer pairs. The fine-tuned model therefore never
saw, in training, the input format it faces in the RAG cells: a question
preceded by four retrieved passages and an instruction to cite them. On
5 of 240 such inputs, the model collapses into a \emph{citation-marker
loop} --- an output consisting of citation markers repeated to the token
limit (``[1] [4] [2] [3] [1] [4] [2] [3] $\ldots$'') --- scoring zero on
questions that the same model, without retrieval, answers correctly. The
failure is diagnostic rather than random: the citation markers are the
most salient novel tokens in the unseen format, and a model trained to
continue the pattern of its prompt continues the markers. We stress that
this is an account of a \emph{training-data omission}, not of retrieval
hurting a well-matched model; Section~\ref{sec:future} describes the
fix (training examples in the exact inference format).

\subsection{Mechanism 2: repetition inherited through self-distillation}

We define \emph{heavy repetition} operationally as a 30-character block
repeated at least $5\times$ within one response. It appears in 35 of 240
DAPT+SFT responses, against 10 of 240 for the base model; the retrieval
cells show 7 of 240 (base+RAG) and 22 of 240 (DAPT+SFT+RAG). Because the
SFT teacher \emph{was} the DAPT model (Section~\ref{sec:sft-data}), the
natural reading is that the teacher's repetition habit was distilled into
the student along with its knowledge: self-distillation transmits
behavior, including degenerate behavior, and a grounding filter on
content words does not screen for it. Repetition is also the more
tractable of the two failure modes --- it is decode-sensitive, and the
demo's repetition penalty of 1.05 (Section~\ref{sec:rag}) largely
suppresses it --- but under the benchmark's unaided greedy decoding it is
part of what is being measured, and we report it rather than decode it
away.

\subsection{Case studies}

The two mechanisms are visible in their purest form on trap items, where
the correct behavior --- rejecting the premise --- is unambiguous.

\textbf{G13 --- ``HA solutions are simple Newtonian fluids, so their
viscosity is a single constant, correct?''} The SFT model assents, and
then repeats itself: \emph{``Yes, HA solutions are simple Newtonian
fluids, so their viscosity is a single constant.''} The same model with
retrieval corrects the premise precisely: \emph{``No, HA solutions are
not simple Newtonian fluids. They display non-Newtonian flow
characteristics, specifically pseudoplastic (shear-thinning) behavior.
This means that their viscosity decreases as shear rate increases.''}

\textbf{G09 --- ``For joint injections, the highest-MW HA is always the
best choice, correct?''} The SFT model correctly rejects the absolute:
\emph{``No, the highest-MW HA is not always the best choice for joint
injections. The choice of HA depends on the specific condition being
treated\ldots''} The SFT+RAG response to the same question degenerates
entirely into a citation-marker loop (``[1] [4] [2] [3] [1] [4] [2] [3]
$\ldots$'', to the token limit). Retrieval, which rescues G13, destroys
G09 --- on a question the parametric model gets right. This is the
interference of Section~\ref{sec:experiments} at the level of a single
response.

\begin{table}[t]
\centering
\begin{tabular}{lcccc}
\toprule
Trap & Base & SFT & Base+RAG & SFT+RAG \\
\midrule
G01 & wrong   & wrong   & correct & correct \\
G02 & wrong   & wrong   & correct & wrong \\
G03 & wrong   & wrong   & correct & partial \\
G04 & partial & correct & correct & correct \\
G05 & partial & correct & correct & correct \\
G06 & partial & correct & correct & wrong \\
G07 & partial & correct & partial & partial \\
G08 & partial & correct & correct & correct \\
G09 & wrong   & correct & partial & wrong \\
G10 & wrong   & ungraded & partial & partial \\
G11 & wrong   & wrong   & correct & correct \\
G12 & correct & ungraded & correct & wrong \\
G13 & wrong   & wrong   & correct & correct \\
G14 & wrong   & wrong   & correct & correct \\
G15 & partial & partial & partial & correct \\
\bottomrule
\end{tabular}
\caption{Per-question verdicts on the 15 trap items. Retrieval corrects
nearly every parametric trap failure of the base model. The SFT+RAG cell
reintroduces failures on items base+RAG answers correctly (G02, G06, G09,
G12) through the interference mechanism of
Section~\ref{sec:analysis}: the trained model's failures with retrieval
are new failures, not residual ones.}
\label{tab:traps}
\end{table}

\section{Limitations}
\label{sec:limitations}

\textbf{The judge is one of the contestants.} Open and trap items are
graded by the SFT model itself (Section~\ref{sec:grading}). Anchoring to
a fixed, evidence-checked reference key --- rather than free pairwise
preference --- and the judge-free auto-graded control (on which the
combined cell ties base+RAG) bound the plausible distortion, but do not
eliminate it. We accordingly read cross-cell gaps below $\approx$0.02 as
noise, and an external judge is the obvious upgrade.

\textbf{Residual key uncertainty.} Verification settled 212 of 240 keys;
28 (11.7\%) remain uncertain drafts, and scores on those items carry key
risk. The same process carries a warning in the other direction:
verification \emph{changed} 22 of 240 keys (9\%) across two rounds. A
benchmark built from unchecked model-written keys would have carried that
error rate silently into every score computed against it --- including,
in our case, into the scores of the very model that wrote the drafts.

\textbf{Small cells.} The multiple-choice column rests on $n=10$: the
SFT+RAG dip there (0.700 vs.\ 0.900) is exactly two questions, and we draw
no conclusion from it. Trap ($n=15$) verdicts are similarly coarse, which
is why we report them per question (Table~\ref{tab:traps}) rather than
only as a mean.

\textbf{The grounding filter checks attestation, not truth.}
Self-generated SFT data can encode teacher errors despite the filter of
Section~\ref{sec:sft-data}: a content-word overlap of $\geq 0.55$
establishes that an answer is about a real passage, not that it states
the passage correctly. The repetition inheritance of
Section~\ref{sec:analysis} is one measured consequence of training on
self-generated data; factual contamination of the same kind cannot be
excluded by our protocol and would require expert audit of the training
set to quantify.

\textbf{Decoding and seed.} All results use a single run with greedy
decoding. Repetition behavior in particular is decode-sensitive (the demo
suppresses it with a repetition penalty that the benchmark deliberately
withholds), so the absolute size of Mechanism 2 should be read as a
property of model-plus-decoding, not of the model alone.

\textbf{Scale and scope.} This is one 8B model in one narrow domain. The
transferable claim is the pipeline and its failure analysis, not any
absolute score, and the interference result in particular should be
re-checked at other scales and in other domains before being treated as
general.

\textbf{Replicability.} Nothing in the pipeline is HA-specific beyond the
corpus query: harvest $\rightarrow$ clean $\rightarrow$ DAPT $\rightarrow$
self-generated grounded SFT $\rightarrow$ evidence-checked bench
$\rightarrow$ four-cell ablation, at a total cloud compute cost of under
USD 50 (single-GPU rental; all other services on free tiers). For small
fields whose literature is open-access, we believe this is a realistic
template for a first trustworthy domain model --- and the interference
result is a warning worth having before, not after, combining SFT with
retrieval.

\section{Future Work}
\label{sec:future}

The failure analysis of Section~\ref{sec:analysis} points to a concrete
remediation path, with one fix per mechanism. For the format mismatch:
\emph{retrieval-formatted SFT data} --- training examples in the exact
inference format (question + retrieved excerpts + cited answer), so the
fine-tuned model learns to use context instead of collapsing on it; we
predict this, alone, should move the combined cell above base+RAG, and
the prediction is directly testable with the released pipeline. For the
repetition inheritance: \emph{repetition filtering} of teacher outputs
before SFT, screening generated pairs with the same operational
definition used to measure the failure, so the habit is not distilled in
the first place; repetition-aware decoding remains available as an
inference-side mitigation. On the evaluation side, re-grading with an
external judge and settling the 28 remaining uncertain keys would tighten
every comparison in Table~\ref{tab:main}.

\section{Data and Code Availability}
\label{sec:availability}

\begin{itemize}
\item Code, corpus manifests, benchmark (v3 keys with per-question
verification status), and full result files:
\url{https://github.com/liukunli001/hyaluronan-expert-llm}
\item Weights: \url{https://huggingface.co/liukunli/ha-expert-qwen3-8b}
(SFT model) and
\url{https://huggingface.co/liukunli/ha-expert-qwen3-8b-dapt} (DAPT
checkpoint)
\item Retrieval index:
\url{https://huggingface.co/datasets/liukunli/ha-expert-index}
\item Interactive demo (chat + cited retrieval):
\url{https://huggingface.co/spaces/liukunli/ha-expert-demo}
\item Corpus release: GitHub Release \texttt{v0.1-data}.
\end{itemize}

\section*{Acknowledgements}

This is an independent research project; the author thanks the open-access
publishers whose PMC deposits made the corpus possible, and the Qwen,
Hugging Face, vLLM, and FAISS communities for the tooling.

\begin{thebibliography}{9}

\bibitem{qwen3}
Qwen Team.
\newblock \emph{Qwen3 Technical Report}.
\newblock Alibaba Group, 2025.

\bibitem{transformers}
Wolf, T., et al.
\newblock Transformers: State-of-the-Art Natural Language Processing.
\newblock In \emph{Proceedings of EMNLP 2020: System Demonstrations}.

\bibitem{trl}
von Werra, L., et al.
\newblock \emph{TRL: Transformer Reinforcement Learning}.
\newblock Hugging Face library, 2023.

\bibitem{neftune}
Jain, N., et al.
\newblock NEFTune: Noisy Embeddings Improve Instruction Finetuning.
\newblock \emph{arXiv:2310.05914}, 2023.

\bibitem{fineweb}
Penedo, G., et al.
\newblock The FineWeb Datasets: Decanting the Web for the Finest Text Data
at Scale.
\newblock In \emph{Advances in Neural Information Processing Systems
(NeurIPS)}, 2024.
\newblock (FineWeb-Edu subset used for replay.)

\bibitem{vllm}
Kwon, W., et al.
\newblock Efficient Memory Management for Large Language Model Serving
with PagedAttention.
\newblock In \emph{Proceedings of SOSP}, 2023.

\bibitem{faiss}
Johnson, J., Douze, M., and J\'egou, H.
\newblock Billion-Scale Similarity Search with GPUs.
\newblock \emph{IEEE Transactions on Big Data}, 7(3):535--547, 2021.

\end{thebibliography}

\end{document}
